\chapter{Variations of Hodge Structure}\label{vhs}
\chapterstatus{complete}

\section{Introduction}\label{vhs:section-introduction}

In a smooth projective family $f \colon \mathcal{X} \to S$, the fibres are diffeomorphic (Ehresmann's theorem, Chapter~\ref{deformations}), so their cohomology groups form a local system with a flat
connection, the Gauss--Manin connection. The Hodge filtrations of the fibres vary holomorphically, and their derivatives move each $F^p$ at most into $F^{p-1}$
(Griffiths transversality). The resulting object, a \emph{variation of Hodge structure}, is recorded by the \emph{period map} to a period domain. Variations
are the natural setting for asking how Hodge classes behave in families. The locus where a class stays of type $(p,p)$ is an analytic subvariety, and by the theorem
of Cattani, Deligne and Kaplan it is algebraic, as the Hodge conjecture predicts. References: \cite[Chapters 9 and 10]{voisin-hodge-1}, \cite{voisin-hodge-2},
\cite{carlson-muller-stach-peters}, \cite{griffiths-periods}.

\noindent\textbf{Prerequisites.} Chapter~\ref{hodge-structures} (Hodge Structures), Chapter~\ref{connections} (Connections and Curvature) and Chapter~\ref{gaga}
(Analytic Spaces and GAGA).

Throughout, $f \colon \mathcal{X} \to S$ is a \emph{smooth projective family}: a proper holomorphic submersion of complex manifolds with connected base $S$, whose fibres
$X_t = f^{-1}(t)$ are projective (for instance the analytification of a smooth projective morphism of varieties).

\section{The Gauss--Manin connection}\label{vhs:section-gauss-manin}

\begin{definition}\label{vhs:definition-gauss-manin}
By Ehresmann's theorem (Theorem~\ref{deformations:theorem-ehresmann}), $\mathcal{H}^k_\ZZ = R^kf_*\ZZ$ is a local system on $S$ with fibres $H^k(X_t; \ZZ)$. The holomorphic vector bundle $\mathcal{H}^k = \mathcal{H}^k_\ZZ \otimes_\ZZ\mathcal{O}_S$ carries the
flat connection $\nabla$ whose flat sections are the local sections of $\mathcal{H}^k_\CC$ (Theorem~\ref{connections:theorem-flat-local-systems}), the \emph{Gauss--Manin connection}. Transport
along a loop gives the \emph{monodromy representation} $\pi_1(S, t) \to \mathrm{GL}(H^k(X_t; \ZZ))$.
\end{definition}

\begin{theorem}\label{vhs:theorem-hodge-bundles}
In a smooth projective family:
\begin{enumerate}
\item the Hodge numbers $h^{p,q}(X_t)$ are independent of $t$;
\item the Hodge filtrations $F^pH^k(X_t; \CC)$ form holomorphic subbundles $\mathcal{F}^p \subseteq \mathcal{H}^k$ (the \emph{Hodge bundles});
\item (\emph{Griffiths transversality}) $\nabla\mathcal{F}^p \subseteq \mathcal{F}^{p-1} \otimes \Omega^1_S$.
\end{enumerate}
\end{theorem}

\begin{proof}[Sketch of proof]
(1) The numbers $h^{p,q}(X_t) = \dim H^q(X_t, \Omega^p_{X_t})$ are upper semicontinuous in $t$: this is Theorem~\ref{schemes:theorem-semicontinuity} for an algebraic family, and
Kodaira--Spencer's semicontinuity theorem in the analytic case \cite[Chapter 9]{voisin-hodge-1}. Their sum over $p + q = k$ is the Betti number $b_k$, constant by Ehresmann's theorem. A sum of
upper semicontinuous integer functions that is constant forces each to be locally constant, hence constant as $S$ is connected. (2) With constant Hodge numbers, the Hodge
filtration of $X_t$ can be computed from harmonic forms for a family of K\"ahler metrics, and depends smoothly on $t$. Holomorphicity follows from (3) applied in the $(0,1)$-direction:
differentiating a family of classes of type $\geq p$ in an antiholomorphic direction does not leave $F^p$ \cite[Chapter 10]{voisin-hodge-1}. (3) Trivialise the family smoothly
(Theorem~\ref{deformations:theorem-ehresmann}) and represent a local section of $\mathcal{F}^p$ by closed forms $\alpha_t$ on the fixed manifold $X_0$ that are of type $\geq p$ for the complex structure
$J_t$. Differentiating in a holomorphic direction $v$ at $t = 0$, $\nabla_v[\alpha] = [\partial_v\alpha_t|_{t=0}]$. A form of type $\geq p$ for $J_t$ is, in terms of $J_0$, a form of type $\geq p$ plus terms
of order $t$ of type $\geq p - 1$, because the $(1,0)$-forms for $J_t$ differ from those for $J_0$ by first-order terms of type $(0,1)$. So $\partial_v\alpha_t|_{t=0}$ has type $\geq p - 1$.
\end{proof}

\begin{definition}\label{vhs:definition-vhs}
A \emph{polarised variation of Hodge structure} of weight $k$ on a complex manifold $S$ is a local system $\mathcal{H}_\ZZ$ of finitely generated abelian groups, a flat bilinear form $Q$, and a
filtration of $\mathcal{H} = \mathcal{H}_\ZZ \otimes \mathcal{O}_S$ by holomorphic subbundles $\mathcal{F}^p$ such that each fibre is a Hodge structure of weight $k$ polarised by $Q$, and
$\nabla\mathcal{F}^p \subseteq \mathcal{F}^{p-1} \otimes \Omega^1_S$. By Theorem~\ref{vhs:theorem-hodge-bundles} and Theorem~\ref{hodge-structures:theorem-cohomology-polarised}, the primitive
cohomology of a smooth projective family, with the polarisation given by a relatively ample class, is one.
\end{definition}

\section{Period domains and period maps}\label{vhs:section-period-domains}

\begin{definition}\label{vhs:definition-period-domain}
Fix a lattice $H_\ZZ$, a bilinear form $Q$ with $Q(y, x) = (-1)^kQ(x, y)$, and Hodge numbers $h^{p,q}$. The \emph{period domain} $D$ is the set of filtrations $F^\bullet$ of $H_\CC$ with
$\dim F^p = \sum_{r \geq p}h^{r,k-r}$ that define Hodge structures of weight $k$ polarised by $Q$. It is an open subset, in the analytic topology, of the \emph{compact dual}
$\check D = \{F^\bullet : Q(F^p, F^{k-p+1}) = 0\}$, a closed subvariety of a flag variety. The real group $G = \mathrm{Aut}(H_\RR, Q)$ acts transitively on $D$ with compact isotropy groups.
\end{definition}

\begin{example}\label{vhs:example-period-domains}
\begin{enumerate}
\item Weight $1$, $h^{1,0} = h^{0,1} = g$, $Q$ symplectic: $D$ is the Siegel upper half space $\mathfrak{H}_g$ (Proposition~\ref{algebraic-groups:proposition-period-matrix}), with
$G = \mathrm{Sp}_{2g}(\RR)$, a Hermitian symmetric domain. For $g = 1$ it is the upper half plane.
\item Weight $2$ with $h^{2,0} = 1$, the case of K3 surfaces ($h^{1,1} = 20$): a Hodge structure is determined by the line $F^2 = H^{2,0} = \CC\omega$, and $D = \{[\omega] \in \PP(H_\CC) : Q(\omega, \omega) = 0,\ Q(\omega, \bar\omega) > 0\}$,
an open subset of a quadric, with $G = \mathrm{SO}(3, 19)$. This is a Hermitian symmetric domain of type IV.
\item For weight $\geq 2$ with $h^{2,0} \geq 2$, or weight $3$ with $h^{3,0} \neq 0$ and $h^{2,1} \neq 0$ (the quintic threefold), $D$ is not Hermitian symmetric, and Griffiths transversality is a
genuine restriction: period maps are tangent to a distribution that is not integrable.
\end{enumerate}
\end{example}

\begin{example}\label{vhs:example-upper-half-plane}
Over the upper half plane $\mathfrak{H}_1 = \{\tau : \operatorname{Im}\tau > 0\}$ consider the family of elliptic curves $E_\tau = \CC/(\ZZ + \ZZ\tau)$, the quotient of $\CC \times \mathfrak{H}_1$ by the action of $\ZZ^2$
given by $(m, n) \cdot (z, \tau) = (z + m + n\tau, \tau)$. The loops $a$, $b$ given by the segments from $0$ to $1$ and from $0$ to $\tau$ form a basis of $H_1(E_\tau; \ZZ)$ that is flat, since the family is
trivialised smoothly by $(x + y\tau, \tau) \mapsto (x, y)$, $x, y \in \RR/\ZZ$. The holomorphic form $dz$ has periods $\int_adz = 1$ and $\int_bdz = \tau$. So in the flat coordinates given by the dual basis of
$H^1$, the line $F^1 = H^{1,0}(E_\tau)$ is spanned by $(1, \tau)$, and the period map $\mathfrak{H}_1 \to \mathfrak{H}_1$ of Example~\ref{vhs:example-period-domains}(1) is the identity. This is the
\emph{tautological} variation: every polarised variation of weight one and rank two is locally the pullback of it by its period map.
\end{example}

\begin{definition}\label{vhs:definition-period-map}
Given a polarised variation on $S$ and a base point, the \emph{period map} $\mathcal{P} \colon \tilde S \to D$ from the universal cover sends $\tilde t$ to the Hodge filtration of the fibre, transported to
$H_\CC = \mathcal{H}_{t_0}$ by the flat connection. It descends to $S \to \Gamma\backslash D$, where $\Gamma$ is the image of the monodromy representation.
\end{definition}

\begin{proposition}\label{vhs:proposition-period-map}
The period map is holomorphic and horizontal: its differential maps $T_tS$ into $\bigoplus_p\Hom(F^p/F^{p+1}, F^{p-1}/F^p) \subseteq T_{\mathcal{P}(t)}\check D$. For a family of projective
manifolds its differential at $t$ is the composite of the Kodaira--Spencer map $T_tS \to H^1(X_t, T_{X_t})$ with the contraction and cup product
$H^1(X_t, T) \to \bigoplus_p\Hom(H^q(\Omega^p), H^{q+1}(\Omega^{p-1}))$.
\end{proposition}

This is Theorem~\ref{vhs:theorem-hodge-bundles}(2), (3) restated, together with Griffiths' computation of the differential \cite[Chapter 10]{voisin-hodge-1}. The second description is
the \emph{infinitesimal variation of Hodge structure}, the main tool in the proof of the Noether--Lefschetz theorem (Theorem~\ref{hodge-conjecture:theorem-noether-lefschetz}).

\begin{example}\label{vhs:example-legendre}
The \emph{Legendre family} $E_\lambda \colon y^2 = x(x - 1)(x - \lambda)$, $\lambda \in \CC \setminus \{0, 1\}$, is a family of elliptic curves. With the holomorphic form $\omega_\lambda = dx/y$
and a symplectic basis $a$, $b$ of $H_1(E_\lambda; \ZZ)$ transported along paths, the period map is $\lambda \mapsto \tau(\lambda) = \int_b\omega_\lambda/\int_a\omega_\lambda \in \mathfrak{H}_1$. The periods
$\int_\gamma\omega_\lambda$ satisfy the Picard--Fuchs equation $\lambda(1 - \lambda)\varpi'' + (1 - 2\lambda)\varpi' - \frac14\varpi = 0$, a differential equation expressing the Gauss--Manin connection.
The monodromy is a subgroup of finite index of $\mathrm{SL}_2(\ZZ)$, the congruence subgroup $\Gamma(2)$.
\end{example}

\begin{proposition}\label{vhs:proposition-k3-transversality}
Let $H_\ZZ$, $Q$ be as in Definition~\ref{vhs:definition-period-domain}, of weight $2$ with $h^{2,0} = 1$. Then for a polarised Hodge structure the filtration is determined by $F^2$, namely
$F^1 = (F^2)^\perp$, and every holomorphic family of such polarised Hodge structures over a complex manifold, with $\mathcal{F}^p$ holomorphic subbundles of a flat bundle $\mathcal{H}$ preserving $Q$ and
$H_\ZZ$, satisfies Griffiths transversality automatically.
\end{proposition}

\begin{proof}
By condition (2) of Definition~\ref{hodge-structures:definition-polarisation}, $Q(H^{2,0}, H^{2,0} \oplus H^{1,1}) = 0$, so $F^1 = H^{2,0} \oplus H^{1,1} \subseteq (F^2)^\perp$; both have dimension
$\dim H - 1$ ($Q$ is nondegenerate), so they are equal. Transversality asks $\nabla\mathcal{F}^2 \subseteq \mathcal{F}^1 \otimes \Omega^1$ and $\nabla\mathcal{F}^1 \subseteq \mathcal{F}^0 \otimes \Omega^1$; the second
is empty since $\mathcal{F}^0 = \mathcal{H}$. For the first, let $\omega$ be a local holomorphic generator of $\mathcal{F}^2$. Then $Q(\omega, \omega) = 0$, and since $Q$ is flat, differentiating along a
vector field $v$ gives $2Q(\nabla_v\omega, \omega) = 0$, so $\nabla_v\omega \in \omega^\perp = \mathcal{F}^1$.
\end{proof}

\begin{counterexample}\label{vhs:counterexample-not-horizontal}
Without such a special shape, transversality is a genuine condition, and holomorphic families of Hodge structures need not be variations of Hodge structure. Let $H_\ZZ = \ZZ^2$ and consider
Hodge structures of weight $2$ with $h^{2,0} = h^{0,2} = 1$ and $h^{1,1} = 0$ (not polarised). Such a structure is a line $F^2 = H^{2,0} = \CC\omega \subseteq \CC^2$ with $\omega$ and $\bar\omega$
independent, so the Hodge structures form $D' = \PP^1(\CC) \setminus \PP^1(\RR)$, two copies of the upper half plane. The tautological family over $D'$, with $\mathcal{F}^2 = \mathcal{F}^1$ the tautological
line in the trivial flat bundle $D' \times \CC^2$, is a holomorphic family of Hodge structures. Transversality would require $\nabla\mathcal{F}^2 \subseteq \mathcal{F}^1 = \mathcal{F}^2$, so that $\mathcal{F}^2$
would be a flat, hence locally constant, subbundle; but the tautological line varies. So this family satisfies all the axioms except polarisation and transversality, and by
Theorem~\ref{vhs:theorem-hodge-bundles}(3) it is not the variation of Hodge structure of any smooth projective family.
\end{counterexample}

\begin{corollary}\label{vhs:corollary-constant-period-map}
If the period map of a polarised variation of Hodge structure is constant, its monodromy group $\Gamma$ is finite. For instance, in the Legendre family (Example~\ref{vhs:example-legendre}) the monodromy is infinite,
so the period map is not constant and the curves $E_\lambda$ are not all isomorphic.
\end{corollary}

\begin{proof}
Transport by the flat connection identifies the fibres over the universal cover with a fixed $(H_\ZZ, Q)$, and the period map is equivariant: for $\gamma \in \pi_1(S, t)$ with monodromy $\rho(\gamma)$,
$\mathcal{P}(\gamma\tilde t) = \rho(\gamma)\mathcal{P}(\tilde t)$ (Definition~\ref{vhs:definition-period-map}). If $\mathcal{P}$ is constant with value $F^\bullet$, then every $\rho(\gamma)$ preserves $F^\bullet$, $H_\ZZ$
and $Q$, so $\Gamma$ lies in the finite group of Proposition~\ref{hodge-structures:proposition-automorphisms-finite}. For the Legendre family $\Gamma = \Gamma(2)$ is infinite (it contains
$\begin{pmatrix}1 & 2\\ 0 & 1\end{pmatrix}$, of infinite order). If all $E_\lambda$ were isomorphic, the $\tau(\lambda)$ would all lie in one orbit of $\mathrm{SL}_2(\ZZ)$ (Section~\ref{algebraic-groups:section-abelian-varieties}),
a countable set; a continuous map from the connected space $\tilde S$ to $\mathfrak{H}_1$ with countable image is constant.
\end{proof}

\section{Global invariants and the theorem of the fixed part}\label{vhs:section-fixed-part}

\begin{theorem}[Theorem of the fixed part]\label{vhs:theorem-fixed-part}
Let $f \colon \mathcal{X} \to S$ be a smooth projective morphism of complex varieties, $\bar{\mathcal{X}}$ a smooth compactification of $\mathcal{X}$, and $t \in S$. Then
$H^0(S, R^kf_*\QQ) = H^k(X_t; \QQ)^{\pi_1(S, t)}$ is a sub-Hodge structure of $H^k(X_t; \QQ)$, and it is the image of the restriction map $H^k(\bar{\mathcal{X}}; \QQ) \to H^k(X_t; \QQ)$
(the \emph{global invariant cycle theorem}).
\end{theorem}

This is \cite[Th\'eor\`eme 4.1.1 and Corollaire 4.1.2]{deligne-hodge-2}; the proof uses mixed Hodge theory (Chapter~\ref{mixed-hodge-structures}). A consequence: a Hodge class on one fibre that
is invariant under monodromy is the restriction of a Hodge class on the total space, and so remains a Hodge class on every fibre. Deligne also proved that the monodromy representation
on $H^k(X_t; \QQ)$ is semisimple.

\section{Hodge loci}\label{vhs:section-hodge-loci}

\begin{definition}\label{vhs:definition-hodge-locus}
Let $\mathcal{H}$ be a variation of Hodge structure of weight $2p$ on $S$. For a flat section $\lambda$ of $\mathcal{H}_\ZZ$ over a simply connected open $U \subseteq S$, the \emph{Hodge locus} of $\lambda$
is $\{t \in U : \lambda_t \in F^p\mathcal{H}_t\}$, the set where $\lambda$ is a Hodge class. The \emph{Hodge locus} of the variation is the union, over all local flat sections, of these sets, a subset of
$S$.
\end{definition}

\begin{proposition}\label{vhs:proposition-hodge-locus-analytic}
The Hodge locus of a flat integral section $\lambda$ over $U$ is a closed analytic subset of $U$, locally defined by $\operatorname{rank}\mathcal{H} - \operatorname{rank}\mathcal{F}^p$ holomorphic equations, so each
component has codimension at most $\dim\mathcal{H}/\mathcal{F}^p = h^{p-1,p+1} + h^{p-2,p+2} + \cdots + h^{0,2p}$. For a real class of weight $2p$, lying in $F^p$ is the same as being of type $(p,p)$.
\end{proposition}

\begin{proof}
The composite $U \to \mathcal{H}|_U \to \mathcal{H}/\mathcal{F}^p$, $t \mapsto$ the image of $\lambda_t$, is a holomorphic section of the holomorphic bundle $\mathcal{H}/\mathcal{F}^p$ of rank
$\sum_{r < p}h^{r,2p-r} = \sum_{j \geq 1}h^{p-j,p+j}$. Its zero set is the Hodge locus, cut out locally by as many holomorphic functions as the rank. For real $\lambda$, $\lambda \in F^p$ implies
$\lambda = \bar\lambda \in \overline{F^p}$, so $\lambda \in F^p \cap \overline{F^p} = H^{p,p}$.
\end{proof}

\begin{theorem}[Cattani--Deligne--Kaplan]\label{vhs:theorem-cdk}
Let $\mathcal{H}$ be a polarised variation of Hodge structure of weight $2p$ over a nonsingular complex algebraic variety $S$ (for instance from a smooth projective family). Then the Hodge locus of $\mathcal{H}$,
restricted to integral classes $\lambda$ with $Q(\lambda, \lambda)$ bounded, is a finite union of algebraic subvarieties of $S$, and the whole Hodge locus is a countable union of algebraic subvarieties.
\end{theorem}

This is \cite{cattani-deligne-kaplan}. The Hodge conjecture would imply it: on the Hodge locus of $\lambda$, the class is algebraic on each fibre, and the locus where a given family of cycles exists is
constructible by the theory of Hilbert schemes, with countably many Hilbert schemes involved. The theorem is proved unconditionally with the several-variable $\mathrm{SL}_2$-orbit theorem describing
the asymptotics of the period map near the boundary of $S$. Recent work of Bakker, Klingler and Tsimerman uses o-minimal geometry to reprove and extend it \cite{bakker-klingler-tsimerman}.

\begin{remark}\label{vhs:remark-variational}
Grothendieck's \emph{variational Hodge conjecture} predicts that if a Hodge class $\lambda$ stays of type $(p,p)$ on all fibres of a family, and is algebraic on one fibre $X_0$, then it is algebraic on all fibres.
By Theorem~\ref{vhs:theorem-fixed-part} it would follow from the Hodge conjecture applied to the total space. It is known for $p = 1$ and in some other cases, and it remains open in general
(Chapter~\ref{known-cases}).
\end{remark}

\section{Exercises}\label{vhs:section-exercises}

\begin{exercise}\label{vhs:exercise-constant-hodge-numbers}
Show that in a smooth projective family of curves the genus is constant, and in a smooth family of K3 surfaces the Hodge numbers of $H^2$ are $(1, 20, 1)$ on every fibre.
\end{exercise}

\begin{solution}
Both follow from Theorem~\ref{vhs:theorem-hodge-bundles}(1): the genus is $h^{1,0}$. For K3 surfaces it suffices that one fibre is a K3 surface with these numbers, for instance a quartic
(Exercise~\ref{calabi-yau:exercise-k3-hodge}).
\end{solution}

\begin{exercise}\label{vhs:exercise-transversality-weight-one}
Show that Griffiths transversality is automatic for variations of weight one: every holomorphic family of Hodge filtrations $F^1 \subseteq H_\CC$ is horizontal.
\end{exercise}

\begin{solution}
For weight one the filtration is $0 = F^2 \subseteq F^1 \subseteq F^0 = H_\CC$, and transversality asks $\nabla F^1 \subseteq F^0 \otimes \Omega^1$, which holds trivially.
\end{solution}

\begin{exercise}\label{vhs:exercise-hodge-locus-codimension}
In a family of K3 surfaces with a nonconstant period map, show that the Hodge locus of an integral class $\lambda \in H^2$ that is not generically of type $(1,1)$ has codimension at most $1$.
\end{exercise}

\begin{solution}
By Proposition~\ref{vhs:proposition-hodge-locus-analytic}, the codimension is at most $\dim\mathcal{H}/\mathcal{F}^1 = h^{0,2} = 1$. Concretely, the locus is where $Q(\lambda, \omega_t) = 0$ for a local
holomorphic generator $\omega_t$ of $F^2$, one holomorphic equation.
\end{solution}

\section{Intermediate Jacobians and Normal Functions}\label{vhs:section-intermediate-jacobians}
Cycle classes forget much of the structure of Chow groups: a cycle whose class vanishes can still be very far from zero. Griffiths' intermediate Jacobian
$J^p(X)$, a complex torus built from the odd-degree Hodge structure $H^{2p-1}(X)$, receives the \emph{Abel--Jacobi map} from homologically trivial cycles of
codimension $p$. For $p = 1$ this is the Picard variety and the classical Abel--Jacobi map. Deligne cohomology packages the cycle class and the Abel--Jacobi
invariant into one group. Families of intermediate Jacobians and their \emph{normal functions} give the only serious analytic method for constructing cycles; it
reproves the Lefschetz $(1,1)$-theorem and proves the Hodge conjecture for cubic fourfolds. References: \cite[Chapter 12]{voisin-hodge-1}, \cite{voisin-hodge-2},
\cite{griffiths-periods}, \cite{carlson-muller-stach-peters}.

\noindent\textbf{Prerequisites.} Chapter~\ref{vhs} (Variations of Hodge Structure), Section~\ref{algebraic-groups:section-abelian-varieties} (Abelian Varieties) and Section~\ref{intersection-theory:section-cycle-class}
(The Cycle Class Map).

Throughout, $X$ is a compact K\"ahler manifold of dimension $n$, usually projective, and $1 \leq p \leq n$.

\section{Intermediate Jacobians}\label{vhs:section-definition}

\begin{definition}\label{vhs:definition-intermediate-jacobian}
The \emph{$p$-th intermediate Jacobian} of $X$ is
\[
J^p(X) = \frac{H^{2p-1}(X; \CC)}{F^pH^{2p-1}(X; \CC) + H^{2p-1}(X; \ZZ)} ,
\]
where $H^{2p-1}(X; \ZZ)$ means its image in $H^{2p-1}(X; \CC)$.
\end{definition}

\begin{proposition}\label{vhs:proposition-torus}
$J^p(X)$ is a compact complex torus of dimension $\frac12b_{2p-1}(X)$. Via Poincar\'e duality,
\[
J^p(X) \cong \frac{F^{n-p+1}H^{2n-2p+1}(X; \CC)^\vee}{H_{2n-2p+1}(X; \ZZ)} ,
\]
where a homology class acts by integration. $J^1(X) = \Pic^0(X)$, and $J^n(X)$ is the \emph{Albanese torus} $\mathrm{Alb}(X) = H^0(X, \Omega^1)^\vee/H_1(X; \ZZ)$.
\end{proposition}

\begin{proof}
$H^{2p-1}$ has odd weight $2p - 1$, so $H_\CC = F^p \oplus \overline{F^{(2p-1)-p+1}} = F^p \oplus \overline{F^p}$ (Proposition~\ref{hodge-structures:proposition-filtration}). Hence the real subspace
$H^{2p-1}(X; \RR)$ meets $F^p$ trivially and maps isomorphically onto $H_\CC/F^p$ (dimensions: $b_{2p-1} = 2\dim_\CC H_\CC/F^p$). So the image of $H^{2p-1}(X; \ZZ)$ modulo torsion is a lattice in
the complex vector space $H_\CC/F^p$, of complex dimension $\frac12b_{2p-1}$. For the second description, the Poincar\'e pairing $H^{2p-1} \times H^{2n-2p+1} \to \CC$ pairs $F^p$ with
$F^{n-p+1}$ to zero (types add up to more than $n$ in one index) and is perfect, so $H_\CC/F^p \cong (F^{n-p+1}H^{2n-2p+1})^\vee$, and $H^{2p-1}(X; \ZZ) \cong H_{2n-2p+1}(X; \ZZ)$ modulo torsion.
For $p = 1$, $H^1(X; \CC)/F^1 = H^{0,1} = H^1(X, \mathcal{O})$, giving $\Pic^0(X)$ (Corollary~\ref{hodge-conjecture:corollary-neron-severi}); for $p = n$ the second description gives
$H^0(X, \Omega^1)^\vee/H_1(X; \ZZ)$.
\end{proof}

\begin{proposition}\label{vhs:proposition-functoriality}
A holomorphic map $f \colon X \to Y$ of compact K\"ahler manifolds induces holomorphic group homomorphisms $f^* \colon J^p(Y) \to J^p(X)$, with $(g \circ f)^* = f^* \circ g^*$ and $\id^* = \id$. If
$H^{2p-1}(X; \QQ) = 0$, then $J^p(X) = 0$.
\end{proposition}

\begin{proof}
The pullback $f^* \colon H^{2p-1}(Y; \CC) \to H^{2p-1}(X; \CC)$ maps $H^{r,s}(Y)$ into $H^{r,s}(X)$ (Proposition~\ref{kahler:proposition-functoriality}), hence $F^pH^{2p-1}(Y)$ into
$F^pH^{2p-1}(X)$, and it maps the image of $H^{2p-1}(Y; \ZZ)$ into the image of $H^{2p-1}(X; \ZZ)$, being induced by the pullback of integral cohomology. So it induces a $\CC$-linear map
$H^{2p-1}(Y; \CC)/F^p \to H^{2p-1}(X; \CC)/F^p$ carrying lattice into lattice, hence a holomorphic homomorphism of the quotient tori; functoriality holds already on cohomology. The last statement is
the dimension formula of Proposition~\ref{vhs:proposition-torus}, since $b_{2p-1}(X) = \dim_\QQ H^{2p-1}(X; \QQ)$.
\end{proof}

\begin{example}\label{vhs:example-albanese-torus}
Let $T = \CC^g/\Lambda$ be a complex torus and $x_0 = 0$. The holomorphic $1$-forms on $T$ are the constant forms $\sum_ja_j\,dz_j$ (Example~\ref{complex-manifolds:example-hodge-numbers} and
Lemma~\ref{complex-manifolds:lemma-harmonic-closed}), so $H^0(T, \Omega^1)^\vee \cong \CC^g$, the functional $\omega \mapsto \int\omega$ being determined by its values on $dz_1, \ldots, dz_g$. For a path $\gamma$ in $T$
from $0$ to $x$, lifted to a path in $\CC^g$ from $0$ to $\tilde x$, $\int_\gamma dz_j = \tilde x_j$; for a closed loop the lift ends at a point of $\Lambda$, and every $\lambda \in \Lambda$ arises from the loop
$t \mapsto t\lambda$. So the period lattice is $\Lambda$, $\mathrm{Alb}(T) = \CC^g/\Lambda = T$, and the Albanese map $x \mapsto \mathrm{AJ}(x - 0)$ of Exercise~\ref{vhs:exercise-aj-points}
is the identity of $T$. Likewise $J^1(T) = \Pic^0(T) = H^{0,1}(T)/H^1(T; \ZZ)$, where $H^{0,1}(T)$ is spanned by the forms $d\bar z_j$ and $H^1(T; \ZZ) = \Hom(\Lambda, \ZZ)$ embeds by sending an $\RR$-linear form on $\CC^g$
with integer values on $\Lambda$ to its $(0,1)$-part; this is the \emph{dual torus} of $T$. For $g = 1$ it is again isomorphic to $T$ (Chapter~\ref{curves}).
\end{example}

\begin{remark}\label{vhs:remark-not-abelian}
The complex structure of $J^p(X)$ is the one for which $F^p$ is killed; it varies holomorphically with $X$ by Griffiths transversality. But $J^p(X)$ is in general not an abelian variety:
the Hodge--Riemann form on $H^{2p-1}$ is definite on each $H^{r,2p-1-r}$ with alternating signs, so it gives a positive definite form on $H_\CC/F^p$ only when the Hodge structure has
\emph{level one}, $H^{2p-1} = H^{p,p-1} \oplus H^{p-1,p}$. Weil's intermediate Jacobian uses a different complex structure that makes it polarised, but it does not vary holomorphically. When
$H^{2p-1}$ has level one, both agree and $J^p(X)$ is a principally polarised abelian variety, for instance for curves, and for $p = 2$ on the cubic threefold
(Example~\ref{vhs:example-cubic-threefold}).
\end{remark}

\section{The Abel--Jacobi map}\label{vhs:section-abel-jacobi}

\begin{definition}\label{vhs:definition-abel-jacobi}
Let $X$ be projective and $Z \in Z^p(X)$ a cycle with $\mathrm{cl}(Z) = 0$ in $H^{2p}(X; \ZZ)$. Then $Z = \partial\Gamma$ for a real $(2n - 2p + 1)$-chain $\Gamma$ (the support of $Z$ has real dimension
$2n - 2p$, and its fundamental class vanishes in $H_{2n-2p}(X; \ZZ)$). The \emph{Abel--Jacobi invariant} is
\[
\mathrm{AJ}(Z) = \Big(\omega \mapsto \int_\Gamma\omega\Big) \in \frac{F^{n-p+1}H^{2n-2p+1}(X; \CC)^\vee}{H_{2n-2p+1}(X; \ZZ)} = J^p(X),
\]
for $\omega$ closed forms of type $\geq n - p + 1$ representing classes in $F^{n-p+1}$.
\end{definition}

\begin{proposition}\label{vhs:proposition-abel-jacobi}
$\mathrm{AJ}(Z)$ is well defined, and $\mathrm{AJ} \colon Z^p(X)_{\mathrm{hom}} \to J^p(X)$ is a homomorphism. It factors through rational equivalence, giving
$\mathrm{AJ} \colon \CH^p(X)_{\mathrm{hom}} \to J^p(X)$. For $p = 1$, and for curves, it is the Abel--Jacobi map of Theorem~\ref{curves:theorem-abel-jacobi}.
\end{proposition}

\begin{proof}[Sketch of proof]
If $\Gamma$ and $\Gamma'$ both bound $Z$, then $\Gamma - \Gamma'$ is a cycle and $\int_{\Gamma-\Gamma'}\omega$ is a period, an element of the image of $H_{2n-2p+1}(X; \ZZ)$. The integral does not depend on the
representative $\omega$ of the class: if $\omega = d\eta$ with $\omega$ of type $\geq n - p + 1$, one can choose $\eta$ of type $\geq n - p + 1$, because the degeneration of the
Hodge-to-de Rham spectral sequence (Corollary~\ref{kahler:corollary-hodge-numbers}) makes $d$ strict for the filtration by type \cite[Chapter 8]{voisin-hodge-1}; and then $\int_\Gamma d\eta = \int_Z\eta = 0$ because $Z$ has complex dimension $n - p$ and $\eta$ has more than $n - p$ holomorphic indices. For
rational equivalence one shows that $t \mapsto \mathrm{AJ}(Z_t - Z_0)$ is a holomorphic map from $\PP^1$ to the torus $J^p(X)$, hence constant, as in the proof of
Theorem~\ref{curves:theorem-abel-jacobi} \cite[Chapter 12]{voisin-hodge-1}.
\end{proof}

\begin{remark}\label{vhs:remark-kernel-abel-jacobi}
For curves, the Abel--Jacobi map $\CH^1(C)_{\mathrm{hom}} \to J^1(C)$ is an isomorphism (Theorem~\ref{curves:theorem-abel-jacobi}). In higher dimension its kernel can be enormous. For $0$-cycles on a
surface $S$ with $h^{2,0}(S) > 0$, Mumford proved that the maps $\mathrm{Sym}^mS \times \mathrm{Sym}^mS \to \CH_0(S)_{\mathrm{hom}}$, $(a, b) \mapsto a - b$, are not surjective for any $m$
\cite{mumford-1968}: the group $\CH_0(S)_{\mathrm{hom}}$ is ``infinite-dimensional'', while its image under the Albanese map $\CH_0(S)_{\mathrm{hom}} \to \mathrm{Alb}(S)$ is a torus of finite dimension. On the other side, Griffiths' theorem (Section~\ref{vhs:section-griffiths}) shows that $\mathrm{AJ}$ can detect cycles which are homologically, but not algebraically,
trivial.
\end{remark}

\section{Deligne cohomology}\label{vhs:section-deligne-cohomology}

\begin{definition}\label{vhs:definition-deligne-complex}
The \emph{Deligne complex} is $\ZZ(p)_{\mathcal{D}} = \big(\ZZ(p) \to \mathcal{O}_X \to \Omega^1_X \to \cdots \to \Omega^{p-1}_X\big)$, with $\ZZ(p) = (2\pi i)^p\ZZ$ in degree $0$, and \emph{Deligne cohomology}
is its hypercohomology $H^k_{\mathcal{D}}(X; \ZZ(p))$. For $p = 1$, $\ZZ(1)_{\mathcal{D}} \simeq \mathcal{O}_X^*[-1]$ by the exponential sequence, so $H^2_{\mathcal{D}}(X; \ZZ(1)) = \Pic(X)$.
\end{definition}

\begin{proposition}\label{vhs:proposition-deligne-sequence}
For $X$ compact K\"ahler there is an exact sequence
\[
0 \to J^p(X) \to H^{2p}_{\mathcal{D}}(X; \ZZ(p)) \to \mathrm{Hdg}^p(X; \ZZ(p)) \to 0 .
\]
For $p = 1$ it is $0 \to \Pic^0(X) \to \Pic(X) \to \mathrm{NS}(X) \to 0$.
\end{proposition}

\begin{proof}
The Deligne complex is the cone, shifted, of $\ZZ(p) \to \Omega^{<p}_X$. So there is a long exact sequence
\[
H^{2p-1}(X; \ZZ(p)) \to \mathbb{H}^{2p-1}(X, \Omega^{<p}) \to H^{2p}_{\mathcal{D}}(X; \ZZ(p)) \to H^{2p}(X; \ZZ(p)) \to \mathbb{H}^{2p}(X, \Omega^{<p}) .
\]
By the holomorphic Poincar\'e lemma and the degeneration of the Hodge-to-de Rham spectral sequence (Corollary~\ref{kahler:corollary-hodge-numbers}),
$\mathbb{H}^k(X, \Omega^{<p}) = H^k(X; \CC)/F^pH^k(X; \CC)$. So the kernel of the last map consists of the integral classes whose image in $H^{2p}(X; \CC)$ lies in $F^p$; a real class in $F^p$ of
weight $2p$ lies in $F^p \cap \overline{F^p} = H^{p,p}$, so these are the integral Hodge classes. The cokernel of the first map is $J^p(X)$ (the twist by $(2\pi i)^p$ only rescales the lattice).
\end{proof}

\begin{example}\label{vhs:example-deligne-low-degree}
\begin{enumerate}
\item For $p = 0$ the Deligne complex is $\ZZ$ in degree $0$, so $H^k_{\mathcal{D}}(X; \ZZ(0)) = H^k(X; \ZZ)$.
\item For $p = 1$, $\ZZ(1)_{\mathcal{D}} \simeq \mathcal{O}_X^*[-1]$ gives $H^k_{\mathcal{D}}(X; \ZZ(1)) = H^{k-1}(X, \mathcal{O}_X^*)$: for compact connected $X$, $H^1_{\mathcal{D}}(X; \ZZ(1)) = \CC^*$ (the
nonzero constants, by Proposition~\ref{complex-manifolds:proposition-no-functions}) and $H^2_{\mathcal{D}}(X; \ZZ(1)) = \Pic(X)$.
\item For $X$ a point and $p \geq 1$, the complex is $\ZZ(p) \to \CC$ in degrees $0$, $1$, so $H^1_{\mathcal{D}}(\mathrm{pt}; \ZZ(p)) = \CC/(2\pi i)^p\ZZ$ and the other groups vanish. For $p = 1$ this is
$\CC/2\pi i\ZZ \cong \CC^*$ via $\exp$, in accordance with (2).
\end{enumerate}
Deligne cohomology is thus a mixture of integral cohomology and of the holomorphic data $\CC/F^p$; it is not a topological invariant, and for $p \geq 1$ its groups are in general neither finitely generated nor
vector spaces.
\end{example}

\begin{remark}\label{vhs:remark-deligne-cycle-class}
There is a \emph{Deligne cycle class} $\mathrm{cl}_{\mathcal{D}} \colon \CH^p(X) \to H^{2p}_{\mathcal{D}}(X; \ZZ(p))$ lifting $\mathrm{cl}$ and restricting to $\mathrm{AJ}$ on homologically trivial cycles
\cite[Chapter 12]{voisin-hodge-1}. In these terms the integral Hodge conjecture asks whether $\mathrm{cl}_{\mathcal{D}}$ followed by the projection to $\mathrm{Hdg}^p(X; \ZZ(p))$ is surjective. For
$p = 1$ both maps are isomorphisms, $\CH^1(X) = \Pic(X) = H^2_{\mathcal{D}}(X; \ZZ(1))$; for $p \geq 2$, $\mathrm{cl}_{\mathcal{D}}$ is neither injective nor surjective in general.
\end{remark}

\section{Griffiths' theorem}\label{vhs:section-griffiths}

\begin{theorem}[Griffiths]\label{vhs:theorem-griffiths}
Let $X \subseteq \PP^4$ be a general quintic threefold. There are homologically trivial $1$-cycles on $X$, differences of lines, that are not algebraically equivalent to zero, even after multiplication
by any nonzero integer. So the Griffiths group $\mathrm{Griff}^2(X) = Z^2(X)_{\mathrm{hom}}/Z^2(X)_{\mathrm{alg}}$ is not torsion.
\end{theorem}

This is \cite{griffiths-1969}. The argument: algebraically trivial cycles map under $\mathrm{AJ}$ to an abelian subvariety of $J^2(X)$ whose tangent space lies in $H^{1,2}$, which for a quintic
threefold is forced to be zero by the variation of Hodge structure. But the Abel--Jacobi images of differences of lines, deformed in the family of quintics, have a nonzero infinitesimal
invariant, so they are not torsion. Clemens later showed that $\mathrm{Griff}^2(X) \otimes \QQ$ is infinite-dimensional for the general quintic \cite{clemens-1983}.

\section{Normal functions}\label{vhs:section-normal-functions}

\begin{definition}\label{vhs:definition-normal-function}
Let $f \colon \mathcal{X} \to S$ be a smooth projective family. The intermediate Jacobians of the fibres form a family of complex tori $\mathcal{J}^p \to S$, the quotient of the holomorphic bundle
$\mathcal{H}^{2p-1}/\mathcal{F}^p$ by the local system $R^{2p-1}f_*\ZZ$. A \emph{normal function} is a holomorphic section $\nu$ of $\mathcal{J}^p$ that is \emph{horizontal}: its local lifts $\tilde\nu$ to
$\mathcal{H}^{2p-1}$ satisfy $\nabla\tilde\nu \in \mathcal{F}^{p-1} \otimes \Omega^1_S$. A family of homologically trivial cycles $Z_t$ gives the normal function $t \mapsto \mathrm{AJ}(Z_t)$.
\end{definition}

\begin{remark}[Poincar\'e and Lefschetz]\label{vhs:remark-poincare-lefschetz}
Lefschetz's original proof of the $(1,1)$-theorem for a surface $S$ used normal functions. Take a Lefschetz pencil of curves $C_t \subseteq S$. An integral class $\lambda$ of type $(1,1)$ restricts to
classes on the $C_t$ that, after subtracting a multiple of the hyperplane class, have degree $0$, and a topological construction produces from $\lambda$ a normal function $\nu$ with
$\nu(t) \in J(C_t) = \Pic^0(C_t)$. By Jacobi inversion (Theorem~\ref{curves:theorem-abel-jacobi}), $\nu(t)$ is the class of a divisor of degree $0$ on $C_t$, and these divisors sweep out a divisor on $S$
whose class is $\lambda$ up to multiples of the hyperplane class (\emph{Poincar\'e's existence theorem}). The same strategy for $p = 2$ fails because the intermediate Jacobians of the
hyperplane sections are not generated by cycles. For cubic fourfolds, where the relevant $J^2$ of the cubic threefold hyperplane sections is covered by Abel--Jacobi images of
lines, Zucker carried it through \cite{zucker-1977}.
\end{remark}

\begin{remark}\label{vhs:remark-green-griffiths}
Green and Griffiths showed that the Hodge conjecture for a primitive Hodge class $\lambda$ of degree $2p$ on $X$ is equivalent to the existence of singularities, over the discriminant locus of a
sufficiently ample linear system, of the normal function attached to $\lambda$ \cite{green-griffiths-2007}. This converts the conjecture into a problem about the boundary behaviour of normal
functions, studied by Brosnan, Pearlstein, Schnell and others.
\end{remark}

\begin{example}\label{vhs:example-cubic-threefold}
For a nonsingular cubic threefold $X \subseteq \PP^4$, $H^3(X)$ has Hodge numbers $h^{3,0} = 0$, $h^{2,1} = 5$, so it has level one and $J^2(X)$ is a principally polarised abelian variety of
dimension $5$. The Abel--Jacobi map from the Fano surface of lines on $X$ to $J^2(X)$ identifies $J^2(X)$ with the Albanese variety of the Fano surface. Clemens and Griffiths showed that
$J^2(X)$ is not a product of Jacobians of curves, and deduced that $X$ is not rational, although it is unirational \cite{clemens-griffiths-1972}.
\end{example}

\section{Exercises on Intermediate Jacobians and Normal Functions}\label{vhs:section-exercises-intermediate-jacobians}

\begin{exercise}\label{vhs:exercise-dimension}
Compute $\dim J^p(X)$ for $X = \PP^n$, for an abelian variety of dimension $g$, and for a quintic threefold with $p = 2$.
\end{exercise}

\begin{solution}
$\PP^n$ has no odd cohomology, so all $J^p = 0$. For an abelian variety, $b_{2p-1} = \binom{2g}{2p-1}$, so $\dim J^p = \frac12\binom{2g}{2p-1}$. For the quintic, $b_3 = 204$
(Proposition~\ref{calabi-yau:proposition-quintic}), so $\dim J^2 = 102$.
\end{solution}

\begin{exercise}\label{vhs:exercise-aj-points}
Show that for $p = n$ the Abel--Jacobi map on $0$-cycles of degree zero is $\sum n_ix_i \mapsto (\omega \mapsto \sum_in_i\int_{x_0}^{x_i}\omega)$, and that it gives the \emph{Albanese map}
$X \to \mathrm{Alb}(X)$, $x \mapsto \mathrm{AJ}(x - x_0)$.
\end{exercise}

\begin{solution}
For $Z = \sum n_i(x_i - x_0)$ with $\sum n_i = 0$, a chain $\Gamma$ with boundary $Z$ is a sum of paths from $x_0$ to $x_i$ with multiplicities $n_i$, and $F^1H^1 = H^{1,0}$ is spanned by holomorphic
$1$-forms (Proposition~\ref{kahler:proposition-holomorphic-forms}). The formula follows from Definition~\ref{vhs:definition-abel-jacobi} with $p = n$.
\end{solution}
